\documentclass[11pt]{article}
\usepackage[margin=1in]{geometry}
\usepackage{booktabs}
\usepackage{array}
\usepackage{amsmath}
\usepackage{url}

% Machine-generated figures (see scripts/diagnostics/make_naics_sic_coverage_brief.R).
% Do NOT hand-edit the \input files; rerun the generator instead.
\newcommand{\covYearMin}{2005}
\newcommand{\covYearMax}{2025}
\newcommand{\covNStates}{56}
\newcommand{\covPermitYears}{139,992}
\newcommand{\covNatNaics}{35.9}
\newcommand{\covNatSic}{98.7}
\newcommand{\covZeroNaics}{16}
\newcommand{\covHiNaics}{10}
\newcommand{\covLoSicList}{ND, VT}
\newcommand{\covSrcFile}{naics\_sic\_coverage\_by\_state\_year\_2026-07-13\_1219.csv}


\title{\vspace{-2.5em}NAICS and SIC Industry-Code Coverage in the NPDES\\
Major-Individual Panel, by State (\covYearMin--\covYearMax)}
\author{}
\date{\today}

\begin{document}
\maketitle
\vspace{-2.5em}

\section*{Scope and data}
This brief measures how often EPA NPDES permits carry an industry classification
code---North American Industry Classification System (\textbf{NAICS}) and the older
Standard Industrial Classification (\textbf{SIC})---and how that coverage varies
across states. The universe is the \textbf{major-individual} permit population
(individual ``NPD'' permits flagged \emph{major} in a given year), reconstructed
by year over \covYearMin--\covYearMax\ with entry and exit allowed. The unit of
observation is a \emph{permit-year}; a permit-year ``has NAICS'' (``has SIC'') if
that permit appears at least once in \texttt{NPDES\_NAICS.csv}
(\texttt{NPDES\_SICS.csv}). Coverage figures pool all \covPermitYears\ permit-years
across the \covNStates\ reporting states and territories.

Two properties of the source matter. First, NAICS/SIC presence is a
\emph{time-invariant permit attribute}: the code files carry no date, so within a
permit, coverage does not change year to year---the cross-state variation below is
structural, not a reporting trend. Second, ``coverage'' means a code is
\emph{on file}; it does not assess whether the code is correct.

\section*{Headline findings}
\begin{itemize}
  \item \textbf{SIC is near-universal: \covNatSic\% nationally}, at or near 100\% in
        almost every state. Only \covLoSicList\ fall below 90\%.
  \item \textbf{NAICS is sparse and sharply state-dependent: \covNatNaics\% nationally.}
        The distribution is bimodal, not centered on the mean: \covZeroNaics\ states
        and territories record NAICS on essentially \emph{none} of their
        major-individual permits, while \covHiNaics\ record it on 95\%+.
  \item \textbf{The gap is administrative, not industrial.} Whether a permit has a
        NAICS code tracks the \emph{permitting authority}, not the facility. Several
        of the largest state programs sit at $\approx$0\% (e.g.\ CA, IL, OH, WI, MD,
        WV; TX at 0.2\%), while neighbors of similar industrial mix (e.g.\ KY, OK, ID
        at 99--100\%) are fully coded.
\end{itemize}

\section*{General state reporting of NAICS and SIC codes}
Industry codes enter ICIS-NPDES at permit issuance and are maintained by the
administering authority---the state environmental agency in delegated (authorized)
states, or the EPA Region otherwise. Because SIC long predates NAICS and was the
required classification when most of these permits and their legacy records were
created, virtually every major-individual permit carries a SIC code: coverage is
\covNatSic\% pooled and effectively saturated in all but a couple of jurisdictions.
NAICS adoption, by contrast, was never uniformly back-filled. A state's ICIS-NPDES
records either systematically include NAICS (yielding near-100\% coverage) or they
do not (yielding near-0\%), producing the two-cluster pattern in
Table~\ref{tab:cov_state} rather than a smooth gradient. The modest national
upward drift over the window---from 33.6\% (2005) to 38.9\% (2025) in
Table~\ref{tab:cov_year}, against flat SIC---reflects newer permits in already-coded
states, not laggard states beginning to report.

\paragraph{Implication for analysis.} For any cross-state industry classification,
use \textbf{SIC}: its coverage is uniform, so it does not confound industry with
state. Pooled NAICS is unusable, and conditioning on NAICS presence would silently
select on state (dropping entire programs such as CA, OH, and IL). This is recorded
in \texttt{docs/data\_quirks.md} as the operative guidance for the panel.

\begin{table}[htbp]
\centering
\scriptsize
\caption{NAICS and SIC coverage of major-individual NPDES permits, pooled
\covYearMin--\covYearMax, by state/territory. $n$ = permit-years with a code on
file; \% = share of that state's permit-years.}\label{tab:cov_state}
\begin{tabular}{l r r r r r}
\toprule
State & Permit-yrs & NAICS $n$ & NAICS \% & SIC $n$ & SIC \% \\
\midrule
AK & 887 & 873 & 98.4 & 866 & 97.6 \\
AL & 3,846 & 2,809 & 73.0 & 3,835 & 99.7 \\
AR & 2,440 & 2,370 & 97.1 & 2,440 & 100.0 \\
AS & 81 & 0 & 0.0 & 81 & 100.0 \\
AZ & 1,495 & 67 & 4.5 & 1,368 & 91.5 \\
CA & 4,909 & 0 & 0.0 & 4,909 & 100.0 \\
CO & 2,608 & 63 & 2.4 & 2,608 & 100.0 \\
CT & 1,904 & 0 & 0.0 & 1,865 & 98.0 \\
DC & 84 & 0 & 0.0 & 84 & 100.0 \\
DE & 328 & 0 & 0.0 & 328 & 100.0 \\
FL & 6,759 & 6,532 & 96.6 & 6,759 & 100.0 \\
GA & 4,175 & 3,937 & 94.3 & 4,101 & 98.2 \\
GU & 97 & 0 & 0.0 & 97 & 100.0 \\
HI & 393 & 0 & 0.0 & 373 & 94.9 \\
IA & 2,761 & 2,714 & 98.3 & 2,714 & 98.3 \\
ID & 820 & 820 & 100.0 & 820 & 100.0 \\
IL & 5,685 & 0 & 0.0 & 5,685 & 100.0 \\
IN & 4,038 & 1,621 & 40.1 & 4,038 & 100.0 \\
KS & 1,065 & 980 & 92.0 & 1,058 & 99.3 \\
KY & 2,874 & 2,864 & 99.7 & 2,874 & 100.0 \\
LA & 4,965 & 24 & 0.5 & 4,965 & 100.0 \\
MA & 2,420 & 63 & 2.6 & 2,420 & 100.0 \\
MD & 1,853 & 0 & 0.0 & 1,853 & 100.0 \\
ME & 1,546 & 396 & 25.6 & 1,546 & 100.0 \\
MI & 3,834 & 3,490 & 91.0 & 3,834 & 100.0 \\
MN & 2,076 & 0 & 0.0 & 1,996 & 96.1 \\
MO & 3,677 & 670 & 18.2 & 3,677 & 100.0 \\
MP & 42 & 0 & 0.0 & 42 & 100.0 \\
MS & 1,894 & 1,752 & 92.5 & 1,894 & 100.0 \\
MT & 709 & 23 & 3.2 & 706 & 99.6 \\
NC & 4,542 & 21 & 0.5 & 4,487 & 98.8 \\
ND & 500 & 385 & 77.0 & 406 & 81.2 \\
NE & 1,045 & 65 & 6.2 & 1,045 & 100.0 \\
NH & 1,010 & 189 & 18.7 & 1,010 & 100.0 \\
NJ & 2,379 & 621 & 26.1 & 2,375 & 99.8 \\
NM & 717 & 84 & 11.7 & 717 & 100.0 \\
NV & 298 & 144 & 48.3 & 298 & 100.0 \\
NY & 6,736 & 1,376 & 20.4 & 6,736 & 100.0 \\
OH & 6,011 & 0 & 0.0 & 6,011 & 100.0 \\
OK & 2,199 & 2,196 & 99.9 & 2,199 & 100.0 \\
OR & 1,518 & 21 & 1.4 & 1,427 & 94.0 \\
PA & 8,072 & 1,417 & 17.6 & 7,796 & 96.6 \\
PR & 1,433 & 0 & 0.0 & 1,433 & 100.0 \\
RI & 499 & 399 & 80.0 & 499 & 100.0 \\
SC & 3,289 & 3,077 & 93.6 & 3,224 & 98.0 \\
SD & 558 & 548 & 98.2 & 558 & 100.0 \\
TN & 3,327 & 3,150 & 94.7 & 3,327 & 100.0 \\
TX & 14,739 & 35 & 0.2 & 14,713 & 99.8 \\
UT & 846 & 841 & 99.4 & 841 & 99.4 \\
VA & 3,043 & 2,706 & 88.9 & 3,043 & 100.0 \\
VI & 121 & 0 & 0.0 & 121 & 100.0 \\
VT & 650 & 629 & 96.8 & 0 & 0.0 \\
WA & 1,424 & 283 & 19.9 & 1,313 & 92.2 \\
WI & 2,468 & 0 & 0.0 & 2,468 & 100.0 \\
WV & 1,905 & 0 & 0.0 & 1,905 & 100.0 \\
WY & 398 & 15 & 3.8 & 382 & 96.0 \\
\midrule
\textbf{All} & 139,992 & 50,270 & 35.9 & 138,170 & 98.7 \\
\bottomrule

\end{tabular}
\end{table}

\begin{table}[h]
\centering
\caption{National coverage by year (pooled across states).}\label{tab:cov_year}
\begin{tabular}{r r r r}
\toprule
Year & Permit-yrs & NAICS \% & SIC \% \\
\midrule
2005 & 5,864 & 33.6 & 99.1 \\
2006 & 6,180 & 33.7 & 98.9 \\
2007 & 6,422 & 33.4 & 98.9 \\
2008 & 6,554 & 33.7 & 98.8 \\
2009 & 6,650 & 33.6 & 98.8 \\
2010 & 6,713 & 33.7 & 98.7 \\
2011 & 6,810 & 34.5 & 98.5 \\
2012 & 6,822 & 34.7 & 98.5 \\
2013 & 6,851 & 35.1 & 98.5 \\
2014 & 6,839 & 35.3 & 98.5 \\
2015 & 6,819 & 35.6 & 98.6 \\
2016 & 6,880 & 36.5 & 98.6 \\
2017 & 6,892 & 36.9 & 98.6 \\
2018 & 6,870 & 37.2 & 98.6 \\
2019 & 6,826 & 37.3 & 98.6 \\
2020 & 6,795 & 37.6 & 98.7 \\
2021 & 6,773 & 37.8 & 98.7 \\
2022 & 6,750 & 37.9 & 98.7 \\
2023 & 6,676 & 37.9 & 98.7 \\
2024 & 6,544 & 38.6 & 98.8 \\
2025 & 6,462 & 38.9 & 98.8 \\
\bottomrule

\end{tabular}
\end{table}

\section*{Caveats}
\begin{itemize}
  \item Figures are \emph{permit-year} shares pooled over the window; a state's
        permit count is its total active major-individual permit-years, not distinct
        permits. Because coverage is time-invariant within a permit, pooling does not
        distort the shares.
  \item Coverage measures \emph{presence} of a code, not its accuracy or its match to
        the facility's true activity.
  \item ``States'' includes DC and U.S.\ territories (\covNStates\ units); the
        smallest have very few permit-years, so their percentages are noisy.
  \item \covLoSicList\ are the only sub-90\% SIC jurisdictions and warrant a look
        before use---atypical against the otherwise saturated SIC pattern.
\end{itemize}

\section*{Reproducibility}
Coverage is computed by
\path{scripts/diagnostics/naics_sic_coverage_by_state_year.R} (writes a
state\,$\times$\,year CSV to \texttt{output/tables/}). This brief's tables are
generated from the newest such CSV (\covSrcFile) by
\path{scripts/diagnostics/make_naics_sic_coverage_brief.R}, which emits the
\texttt{\textbackslash input} fragments and the in-text numbers. Rerun that
generator after any refresh; do not edit the fragments or this PDF by hand.

\end{document}
